% บทที่ 3 วิธีดำเนินงาน
% ย้ายมาจากเค้าโครง ../chapter/04-methodology.tex แล้วแก้ให้ตรงกับที่ทำจริงใน codingpy/
\chapter{วิธีดำเนินงาน}

โครงงานนี้ทำตามกระบวนการวิทยาการข้อมูล 5 ขั้นตอน
คือ การตั้งคำถาม การเก็บรวบรวมข้อมูล การสำรวจและทำความสะอาดข้อมูล การวิเคราะห์เชิงสถิติ และการสรุปผล
บทนี้อธิบายสี่ขั้นแรก ผลของการวิเคราะห์อยู่ในบทที่ 4

\section{คำถามวิจัย}

\begin{subitems}
    \item เวลาที่ใช้กับโซเชียลมีเดียต่อวัน สัมพันธ์กับคะแนนภาวะซึมเศร้า ความวิตกกังวล และความเครียดอย่างไร และแรงเพียงใด
    \item โซเชียลมีเดียสัมพันธ์กับมิติใดของ DASS-21 มากที่สุด
    \item การติดตามข่าวเชิงลบสัมพันธ์กับความวิตกกังวลมากกว่ามิติอื่นหรือไม่
    \item นักศึกษาที่เสพเนื้อหาต่างประเภทกัน มีคะแนนทั้งสามมิติแตกต่างกันหรือไม่
    \item เวลาที่ใช้กับโซเชียลมีเดียสัมพันธ์กับจำนวนชั่วโมงและคุณภาพการนอนอย่างไร
    \item เมื่อควบคุมการนอน การออกกำลังกาย ความกดดันจากการเรียน ผลการเรียน และความเพียงพอทางการเงินแล้ว เวลาที่ใช้กับโซเชียลมีเดียยังอธิบายความแตกต่างของคะแนนแต่ละมิติได้อยู่หรือไม่
    \item ในกลุ่มนักศึกษาที่มีผลการเรียนและความเพียงพอทางการเงินใกล้เคียงกัน ผู้ที่ใช้โซเชียลมีเดียมากกว่ามีคะแนนความเครียดต่างจากผู้ที่ใช้น้อยกว่าหรือไม่
\end{subitems}

\section{การเก็บรวบรวมข้อมูล}

\subsection{เหตุผลที่เก็บข้อมูลเอง}

ชุดข้อมูลสาธารณะในหัวข้อนี้บน Kaggle \cite{singh_student_social_media,singh_social_media_addiction}
มีลักษณะของข้อมูลสังเคราะห์หลายข้อ
หลักทศนิยมกระจายสม่ำเสมอทั้ง 10 หลัก ทั้งที่คนจริงมักตอบเป็นเลขกลม
แพลตฟอร์มผูกกับประเทศแบบตายตัวทุกแถว เช่น LINE ปรากฏเฉพาะในญี่ปุ่น
ตัวแปรที่วัดคนละเรื่องมีสหสัมพันธ์กันสูงถึง 0.80 ถึง 0.95
และไม่มีค่าว่างเลยสักช่อง
ข้อมูลแบบนี้อธิบายสถานการณ์จริงของนักศึกษาไม่ได้
คณะผู้จัดทำจึงเก็บข้อมูลเอง และเก็บเฉพาะตัวแปรที่ต้องใช้ตอบคำถามวิจัย

\subsection{กลุ่มตัวอย่าง}

ผู้ตอบคือนักศึกษาระดับปริญญาตรี มหาวิทยาลัยขอนแก่น
เลือกแบบตามสะดวก (convenience sampling) โดยส่งลิงก์แบบสอบถามในกลุ่มออนไลน์ของคณะและชั้นปีต่าง~ๆ
เค้าโครงตั้งเป้าไว้ 150 ถึง 200 คน แต่ได้คำตอบจริง 121 คน
ระหว่างวันที่ 23 สิงหาคม ถึง 3 กันยายน 2569
และเหลือ 107 คนหลังคัดผู้ที่ไม่ผ่านข้อดักความใส่ใจออก (หัวข้อ~\ref{sec:cleaning})

\subsection{เครื่องมือ}

แบบสอบถามเป็น Google Forms 36 ข้อ ใช้เวลาตอบราว 8 ถึง 9 นาที
ตารางที่~\ref{tab:form} แสดงโครงสร้าง

\begin{table}[H]
\centering
\caption{โครงสร้างแบบสอบถาม}
\label{tab:form}
\begin{tabularx}{\textwidth}{|c|L|l|}
\hline
\textbf{ข้อ} & \textbf{คำถาม} & \textbf{รูปแบบคำตอบ} \\
\hline
-- & คำชี้แจงและการขอความยินยอม & บังคับตอบ \\
\hline
1--4 & อายุ เพศ ชั้นปี คณะ & ตัวเลข/ตัวเลือก \\
\hline
5 & ชั่วโมงใช้โซเชียลมีเดียต่อวัน (จาก Screen Time) & ตัวเลข \\
\hline
6 & แพลตฟอร์มที่ใช้มากที่สุด & ตัวเลือกเดียว \\
\hline
7 & ประเภทเนื้อหาที่ดูมากที่สุด & ตัวเลือกเดียว 8 ข้อ \\
\hline
8 & จำนวนวันใน 7 วันที่ติดตามข่าวหนัก & ตัวเลข 0--7 \\
\hline
9--10 & ชั่วโมงการนอน และคุณภาพการนอน & ตัวเลข / สเกล 1--10 \\
\hline
11--12 & จำนวนวันที่ออกกำลังกาย และความกดดันจากการเรียน & ตัวเลข / สเกล 1--10 \\
\hline
13 & เกรดเฉลี่ยสะสม & ตัวเลข 0.00--4.00 \\
\hline
14 & ความเพียงพอของเงินที่ได้รับต่อเดือน & สเกล 1--10 \\
\hline
15 & ข้อดักความใส่ใจ & ตัวเลือกเดียว \\
\hline
16--36 & แบบวัด DASS-21 & สเกล 0--3 \\
\hline
\end{tabularx}
\end{table}

ตัวเลือกของข้อ 7 มี 8 ข้อ คือ ข่าวสารบ้านเมืองและการเมือง บันเทิง เกมและอีสปอร์ต กีฬา
ความรู้และการพัฒนาตนเอง ชีวิตประจำวันของคนอื่น สินค้าและรีวิว และอื่น~ๆ

ข้อ 5 ไม่ได้ให้ผู้ตอบประมาณจากความจำ
แบบสอบถามบอกให้เปิดดูค่าที่มือถือบันทึกไว้
(iPhone: ตั้งค่า $\rightarrow$ Screen Time \quad Android: ตั้งค่า $\rightarrow$ Digital Wellbeing)
เพื่อลดความคลาดเคลื่อนจากการรายงานตนเอง

แบบสอบถามบังคับตอบเฉพาะข้อที่ขาดไม่ได้ ได้แก่ การยินยอม ข้อ 5 ข้อ 15 และแบบวัด DASS-21 ทั้ง 21 ข้อ
เพราะถ้าขาด DASS-21 ไปข้อเดียวก็คำนวณคะแนนของมิตินั้นไม่ได้
ข้ออื่นไม่บังคับ เหตุผลคือผู้ตอบที่ไม่อยากตอบข้อใดข้อหนึ่งมักปิดแบบสอบถามทิ้งทั้งชุด
หรือเลือกตอบส่งเดชให้ผ่านไป ค่าว่างที่มองเห็นได้จัดการง่ายกว่าคำตอบที่ไม่ตั้งใจ
และการบังคับทุกข้อยังขัดกับหลักที่ว่าผู้ตอบหยุดตอบได้ตลอดเวลา

\subsection{ข้อดักความใส่ใจ}

ข้อ 15 มีข้อความดังนี้

\begin{quote}
``ข้อนี้ไม่ใช่คำถาม แต่เป็นข้อตรวจสอบว่าท่านอ่านข้อความจริง
กรุณาเลือกตัวเลือกที่ 3 \textbf{ไม่ค่อยตรงกับฉัน}''

ตัวเลือก: (1) ตรงกับฉันมาก \quad (2) ค่อนข้างตรงกับฉัน
\quad (3) ไม่ค่อยตรงกับฉัน \quad (4) ไม่ตรงกับฉันเลย
\end{quote}

ข้อความนี้ผ่านการทดลองใช้และปรับแก้มาหนึ่งรอบ
ฉบับแรกมีข้อความว่า ``ตรวจสอบว่าแบบสอบถามที่ผ่านมาตรงกับคุณมากแค่ไหน''
ผู้ทดลองตอบเข้าใจว่าเป็นคำถามที่ต้องตอบตามความรู้สึกจริง ไม่ใช่คำสั่ง
จึงตัดข้อความนั้นออก ขึ้นต้นด้วย ``ข้อนี้ไม่ใช่คำถาม''
และเขียนทั้งหมายเลขและข้อความของตัวเลือกกำกับไว้
ข้อดักที่กำกวมจะคัดผู้ที่ตั้งใจตอบออกไปด้วย ซึ่งแย่กว่าไม่มีข้อดักเลย

คำตอบที่ถูกอยู่ที่ตัวเลือกที่ 3 ไม่ใช่ตัวเลือกแรกหรือสุดท้าย
เพราะคนที่กดโดยไม่อ่านมักกดตำแหน่งเดิมซ้ำ ซึ่งมักเป็นหัวแถวหรือท้ายแถว
สเกล 4 ระดับทำให้ข้อนี้หน้าตาคล้ายแบบวัด DASS-21 ที่ตามมา
และวางไว้ก่อนแบบวัด ไม่แทรกระหว่างข้อ เพื่อไม่ให้กระทบความเชื่อมั่นของแบบวัด

\subsection{แบบวัด DASS-21}

โครงงานเลือก DASS-21 ด้วยเหตุผลสี่ข้อ

\begin{subsubitems}
    \item ให้คะแนนสามมิติแยกกัน จึงบอกได้ว่าโซเชียลมีเดียสัมพันธ์กับมิติใดมากกว่ากัน
    \item มิติความวิตกกังวลตรงกับคำถามเรื่องการเสพข่าวเชิงลบ และมิติความเครียดตรงกับประเด็นในที่มาของโครงงาน
    \item ไม่มีข้อคำถามเกี่ยวกับการทำร้ายตนเอง จึงเหมาะกับการเก็บข้อมูลโดยนักศึกษาที่ไม่มีบุคลากรสุขภาพจิตคอยดูแล
    \item มีฉบับแปลภาษาไทยบนเว็บไซต์ทางการของแบบวัด \cite{webster2008thai}
          และผ่านการตรวจสอบคุณภาพในนักศึกษาไทยแล้ว \cite{wittayapun2023validation}
\end{subsubitems}

แบบสอบถามใช้ข้อความฉบับแปลของ Webster \cite{webster2008thai}
โดยไม่แก้ข้อความ ไม่ตัดข้อ ไม่สลับลำดับ และไม่เปลี่ยนสเกลคำตอบ
แต่ละข้อตอบ 0 (ไม่เคยเกิดขึ้นเลย) ถึง 3 (เกิดขึ้นบ่อยมากหรือเกือบตลอดเวลา)
ตามความรู้สึกในช่วง 1 สัปดาห์ที่ผ่านมา แบบวัด DASS เผยแพร่ให้ใช้ในงานวิจัยได้โดยไม่ต้องขออนุญาต
ข้อแลกคือแบบสอบถามยาวขึ้นจาก 18 ข้อเป็น 36 ข้อ

\subsection{ตัวแปร}

ตารางที่~\ref{tab:variables} แสดงตัวแปรที่ใช้ในการวิเคราะห์ ชื่อตัวแปรตรงกับชื่อคอลัมน์ในไฟล์ข้อมูล

\begin{table}[H]
\centering
\caption{ตัวแปรที่ใช้ในการวิเคราะห์}
\label{tab:variables}
\begin{tabularx}{\textwidth}{|l|L|l|c|}
\hline
\textbf{ตัวแปร} & \textbf{คำอธิบาย} & \textbf{ระดับการวัด} & \textbf{ช่วงค่า} \\
\hline
\multicolumn{4}{|l|}{\textit{ตัวแปรตาม}} \\
\hline
\texttt{dass\_depression} & คะแนนภาวะซึมเศร้า & อันตรภาค & 0--42 \\
\hline
\texttt{dass\_anxiety} & คะแนนความวิตกกังวล & อันตรภาค & 0--42 \\
\hline
\texttt{dass\_stress} & คะแนนความเครียด & อันตรภาค & 0--42 \\
\hline
\multicolumn{4}{|l|}{\textit{ตัวแปรต้นหลัก}} \\
\hline
\texttt{social\_media\_hours} & ชั่วโมงใช้โซเชียลมีเดียต่อวัน อ่านจาก Screen Time & อัตราส่วน & 0--24 \\
\hline
\texttt{news\_days} & จำนวนวันใน 7 วันที่ผ่านมาที่ติดตามข่าวเชิงลบ & อัตราส่วน & 0--7 \\
\hline
\texttt{content\_group} & ประเภทเนื้อหาที่ดูมากที่สุด ยุบเหลือ 4 กลุ่ม & นามบัญญัติ & -- \\
\hline
\multicolumn{4}{|l|}{\textit{ตัวแปรควบคุม}} \\
\hline
\texttt{sleep\_hours} & ชั่วโมงการนอนเฉลี่ยต่อคืน & อัตราส่วน & 0--24 \\
\hline
\texttt{sleep\_quality} & คุณภาพการนอนโดยรวม & เรียงลำดับ & 1--10 \\
\hline
\texttt{exercise\_days} & จำนวนวันที่ออกกำลังกายต่อสัปดาห์ & อัตราส่วน & 0--7 \\
\hline
\texttt{study\_pressure} & ความกดดันจากการเรียน & เรียงลำดับ & 1--10 \\
\hline
\texttt{GPA} & เกรดเฉลี่ยสะสม & อัตราส่วน & 0--4 \\
\hline
\texttt{financial\_sufficiency} & ความเพียงพอของเงินเทียบกับค่าใช้จ่ายของตัวเอง & เรียงลำดับ & 1--10 \\
\hline
\end{tabularx}
\end{table}

DASS-21 เป็นแบบวัดเชิงลบ คะแนนยิ่งสูงแปลว่าอาการยิ่งมาก
ส่วน \texttt{sleep\_quality} และ \texttt{financial\_sufficiency} ยิ่งสูงยิ่งดี
ค่าสัมประสิทธิ์ที่ติดลบของสองตัวนี้จึงแปลว่ายิ่งนอนดีหรือเงินพอ อาการยิ่งน้อย

ตัวแปรควบคุมทั้งหกตัวมีไว้กันไม่ให้สรุปผิดว่าผลมาจากโซเชียลมีเดีย
ทั้งที่จริงมาจากการนอนน้อย ความกดดันจากการเรียน หรือเรื่องเงิน
ผลการเรียนและการเงินเป็นแหล่งความเครียดของนักศึกษา และอาจสัมพันธ์กับเวลาที่ใช้กับโซเชียลมีเดียไปพร้อมกัน
จึงเข้าข่ายเป็นตัวแปรกวน
เรื่องการเงินถามเป็นระดับความเพียงพอเทียบกับค่าใช้จ่ายของผู้ตอบเอง ไม่ได้ถามจำนวนบาท
เพราะเงินเท่ากันมีความหมายต่างกันตามภาระ เช่น คนพักหอกับคนพักบ้าน
และคำถามจำนวนเงินมักถูกปฏิเสธ

\section{การทำความสะอาดข้อมูล}
\label{sec:cleaning}

สคริปต์ \texttt{clean\_data.py} แปลงคำตอบเป็นตัวเลข
ตามกติกาในตารางที่~\ref{tab:cleaning}

\begin{table}[H]
\centering
\caption{กติกาการทำความสะอาดข้อมูล}
\label{tab:cleaning}
\begin{tabularx}{\textwidth}{|L|L|}
\hline
\textbf{ปัญหา} & \textbf{วิธีจัดการ} \\
\hline
คำตอบเป็นข้อความปนตัวเลข เช่น ``8.5 ชั่วโมง'' หรือช่วง ``6-7 ชม.'' & ดึงตัวเลขทุกตัวในช่องแล้วหาค่าเฉลี่ย ช่วง 6--7 จึงเป็น 6.5 \\
\hline
ค่าที่เป็นไปไม่ได้ เช่น นอน 210 ชั่วโมง หรือเกรดเฉลี่ยเกิน 4.00 & กำหนดเป็นค่าว่าง \\
\hline
ไม่มีผลการเรียน หรือตอบ ``-'' ในข้อเกรดเฉลี่ย & กำหนดเป็นค่าว่าง \\
\hline
ผู้ตอบไม่ผ่านข้อดักความใส่ใจ & ตัดออกทั้งแถว \\
\hline
ประเภทเนื้อหาบางตัวเลือกมีผู้ตอบน้อย & ยุบเหลือ 4 กลุ่มที่กำหนดไว้ก่อนเห็นผล คือ ข่าว บันเทิง ชีวิตคนอื่น และอื่น~ๆ \\
\hline
ค่าว่างในข้อที่ไม่บังคับตอบ & ไม่เติมค่า แต่ละการวิเคราะห์ใช้แถวที่มีข้อมูลครบเฉพาะตัวแปรที่ใช้ \\
\hline
\end{tabularx}
\end{table}

การยุบประเภทเนื้อหาใช้กติกาที่กำหนดไว้ในเค้าโครงก่อนเห็นผล
เพื่อไม่ให้เป็นการเลือกยุบกลุ่มตามผลที่อยากได้
ส่วนการถดถอยพหุคูณใช้แถวที่ครบทุกตัวแปรชุดเดียวกันทั้งสามสมการ
เพื่อให้เทียบ $R^2$ และสัมประสิทธิ์ข้ามสมการได้

คะแนน DASS-21 รายมิติคำนวณจากผลรวมของข้อในมิตินั้นคูณ 2 ตามคู่มือ \cite{lovibond1995manual}
ภาวะซึมเศร้าคือข้อ 3, 5, 10, 13, 16, 17, 21
ความวิตกกังวลคือข้อ 2, 4, 7, 9, 15, 19, 20
และความเครียดคือข้อ 1, 6, 8, 11, 12, 14, 18
ส่วนระดับความรุนแรงของแต่ละมิติแบ่งตามเกณฑ์ในตารางที่~\ref{tab:severity}

\begin{table}[H]
\centering
\caption{เกณฑ์แบ่งระดับความรุนแรงของ DASS-21 (คะแนนหลังคูณ 2)}
\label{tab:severity}
\begin{tabular}{|l|c|c|c|}
\hline
\textbf{ระดับ} & \textbf{ซึมเศร้า} & \textbf{วิตกกังวล} & \textbf{เครียด} \\
\hline
ปกติ & 0--9 & 0--7 & 0--14 \\
\hline
เล็กน้อย & 10--13 & 8--9 & 15--18 \\
\hline
ปานกลาง & 14--20 & 10--14 & 19--25 \\
\hline
รุนแรง & 21--27 & 15--19 & 26--33 \\
\hline
รุนแรงมาก & 28 ขึ้นไป & 20 ขึ้นไป & 34 ขึ้นไป \\
\hline
\end{tabular}
\end{table}

ความเชื่อมั่นของแต่ละมิติคำนวณด้วยค่าแอลฟาของครอนบาค (\texttt{reliability.py})
พร้อมช่วงความเชื่อมั่น 95\% แบบ bootstrap 5,000 รอบ
และดูรายข้อว่าข้อใดสัมพันธ์กับข้ออื่นในมิติเดียวกันต่ำ หรือทำให้แอลฟาสูงขึ้นถ้าตัดออก

\section{การตรวจสอบคุณภาพและที่มาของข้อมูล}

ข้อมูลที่เก็บเองถูกเทียบกับชุดข้อมูลสังเคราะห์ด้วยเกณฑ์สามข้อ
ข้อแรกคือการกระจายของหลักทศนิยมในชั่วโมงการใช้งาน คนจริงมักตอบเป็นเลขกลม
ข้อที่สองคือสัดส่วนค่าว่าง แบบสอบถามจริงย่อมมีค่าว่างบ้าง
ข้อที่สามคือขนาดของสหสัมพันธ์ระหว่างตัวแปรต้น ซึ่งไม่ควรสูงถึง 0.80 ถึง 0.95 แบบชุดข้อมูลสังเคราะห์

\section{การวิเคราะห์ข้อมูล}

การวิเคราะห์ทุกข้อทำกับตัวแปรตามสามตัวแยกกัน
ตารางที่~\ref{tab:analysis} สรุปวิธีและสคริปต์แยกตามคำถาม

\begin{table}[H]
\centering
\caption{วิธีวิเคราะห์แยกตามคำถามวิจัย}
\label{tab:analysis}
\begin{tabularx}{\textwidth}{|c|L|l|}
\hline
\textbf{ข้อ} & \textbf{วิธีวิเคราะห์} & \textbf{สคริปต์} \\
\hline
1, 2 & สหสัมพันธ์เพียร์สันระหว่างชั่วโมงใช้งานกับคะแนนทั้งสามมิติ พร้อม $r^2$ ช่วงความเชื่อมั่น 95\% แบบ bootstrap สหสัมพันธ์สเปียร์แมน และสหสัมพันธ์บางส่วน & \texttt{qr2.py}, \texttt{q1\_usage\_dass.py} \\
\hline
3 & วิธีเดียวกับข้อ 1 โดยใช้จำนวนวันที่ติดตามข่าวเชิงลบ & \texttt{qr2.py}, \texttt{q3\_news\_anxiety.py} \\
\hline
4 & Kruskal-Wallis เทียบ 4 กลุ่มเนื้อหาทั้งสามมิติ และ ANOVA สำหรับมิติความวิตกกังวล & \texttt{q4\_content\_type.py}, \texttt{hypothesis\_tests.py} \\
\hline
5 & สหสัมพันธ์ระหว่างชั่วโมงใช้งานกับชั่วโมงและคุณภาพการนอน & \texttt{q5\_interaction\_dass.py} \\
\hline
6 & การถดถอยพหุคูณแบบลำดับขั้น (hierarchical) 3 สมการ & \texttt{q6\_regression.ipynb} \\
\hline
7 & สัมประสิทธิ์ของ \texttt{social\_media\_hours} ในสมการความเครียด และ partial regression plot & \texttt{q6\_regression.ipynb} \\
\hline
เพิ่มเติม & ANOVA เทียบคะแนนตามสภาพการเงิน 3 กลุ่ม และสหสัมพันธ์กับเกรดเฉลี่ย & \texttt{qr7.py}, \texttt{q6\_gpa\_finance.py} \\
\hline
\end{tabularx}
\end{table}

สหสัมพันธ์บางส่วนในข้อ 1 ถึง 3 ควบคุมชั่วโมงการนอน คุณภาพการนอน ความกดดันจากการเรียน
เกรดเฉลี่ย และความเพียงพอทางการเงิน

การถดถอยในข้อ 6 ทำเป็นสองขั้น ขั้นแรกใส่ตัวแปรควบคุมหกตัว
ขั้นที่สองเพิ่มตัวแปรโซเชียลมีเดียสองตัว คือ \texttt{social\_media\_hours} และ \texttt{news\_days}
แล้วดูว่า $R^2$ เพิ่มขึ้นเท่าใด ($\Delta R^2$) และการเพิ่มนั้นมีนัยสำคัญหรือไม่ด้วย F-change
สมการเต็มของแต่ละมิติจึงเป็น

\begin{center}
\texttt{dass\_<มิติ> $\sim$ social\_media\_hours + news\_days + sleep\_hours}\\
\texttt{+ sleep\_quality + exercise\_days + study\_pressure}\\
\texttt{+ GPA + financial\_sufficiency}
\end{center}

\noindent
คะแนน DASS-21 เบ้ขวาและความแปรปรวนของค่าคลาดเคลื่อนอาจไม่คงที่
ค่าความคลาดเคลื่อนมาตรฐานของสัมประสิทธิ์จึงใช้แบบ HC3 ซึ่งทนต่อปัญหานี้
ค่า $p$ ของการทดสอบชุดหลักปรับด้วย Bonferroni โดยคูณ 3 เพราะทดสอบ 3 สมการ
ข้อตกลงเบื้องต้นตรวจด้วย VIF, residual plot, Q-Q plot, การทดสอบ Breusch-Pagan
และการทดสอบ Shapiro-Wilk ของค่าคลาดเคลื่อน

นอกจากคำถามวิจัยทั้งเจ็ดข้อ รายวิชากำหนดให้ทดสอบสมมติฐานสี่แบบ
คือ ค่าเฉลี่ยประชากรเดียว สัดส่วนประชากรเดียว ค่าเฉลี่ยสองประชากร และค่าเฉลี่ยหลายประชากร
ทุกการทดสอบใช้ระดับนัยสำคัญ 0.05 และอยู่ในสคริปต์ \texttt{hypothesis\_tests.py}

เครื่องมือที่ใช้คือ Python จัดการแพ็กเกจด้วย uv
ร่วมกับไลบรารี pandas, numpy, scipy, statsmodels และ matplotlib
ทุกสคริปต์ตั้งค่า seed เป็น 42 ผลที่ได้จึงรันซ้ำแล้วได้ค่าเดิม
